\documentclass[11pt]{article}
\usepackage[margin=1in]{geometry}
\usepackage{amsmath,amssymb,amsfonts,mathtools}
\usepackage{array,booktabs,enumitem,graphicx,xcolor,hyperref}
\usepackage[T1]{fontenc}
\usepackage[utf8]{inputenc}
\title{Discrete and Continuum Approaches to Three-Dimensional Quantum Gravity}
\author{Hirosi Ooguri and Naoki Sasakura}
\date{}
\begin{document}
\maketitle
\begin{abstract}

It is shown that, in the three-dimensional lattice gravity defined by Ponzano and Regge, the space of physical states is isomorphic to the space of gauge-invariant functions on the moduli space of flat \$SU(2)\$ connections over a two-dimensional surface, which gives physical states in the \$ISO(3)\$ Chern-Simons gauge theory.

\end{abstract}

\section{Introduction}

= @big=11.9 @ @Big=16.1 @ @bigg=20.3 @ @Bigg=24.5 @ = @nt =cmr10 scaled 4 =cmbx10 scaled 4 =cmsl10 scaled 4 =cmti10 scaled 4 =cmss12 =cmss9 =cmss8 ="090B ="090C ="090D ="090E ="090F ="0910 ="0911 ="0912 ="0913 ="0914 ="0915 ="0916 ="0917 ="0918 ="091A ="091B ="091C ="091D ="091E ="091F ="0920 ="0921

\section{Method}

V.G.\textasciitilde{}Drinfeld, `` Quantum Groups ,'' in Proc.\textasciitilde{}Intl.\textasciitilde{}Congress of Mathematicians , Berkeley, California, (1986) p.798. V.G.\textasciitilde{}Turaev and O.Y.\textasciitilde{}Viro, `` State Sum Invariants of 3-Manifolds and Quantum math expression -Symbols, '' preprint (1990). J.W.\textasciitilde{}Alexander, Ann. Math. 31 (1930) 292. A.\textasciitilde{}Tsuchiya and Y.\textasciitilde{}Kanie, in Conformal Field Theory and Solvable Lattice Models , Advanced Studies in Pure Mathematics, 16 (1988) 297;

\begin{equation}
\label{eq:compact}
L(\theta) = \sum_{i=1}^{n} \ell(x_i, \theta)
\end{equation}
Equation~\ref{eq:compact} is included to keep the sample paper-like while remaining small.

\section{Conclusion}

math expression where math expression is a dreibein, math expression is a spin connection on math expression , and math expression is a curvature two-form computed from math expression as math expression . We will examine physical states in the Ponzano-Regge model and show that they can be identified with the ones in the continuum field theory . There is a disturbing factor of math expression in front of the Einstein-Hilbert action math expression in the exponent of the integrand. We shall discuss on this issue at the end of this letter. Now let us discuss the regularization of the sum over coloring. The method adopted by Ponzano and Regge in their original paper was simply to cut-off the sum by math expression and rescale it by multiplying a factor math expression per each vertex. math expression Z(L) = \& \_ colorings (j\_e L) \textasciitilde{}\textasciitilde{}\textasciitilde{} \_ vertices 1 (L) \textasciitilde{} \_ e:edges (-1)\textasciicircum{} 2j\_e (2j\_e+1) \& \textasciitilde{}\textasciitilde{}\textasciitilde{}\textasciitilde{}\textasciitilde{}\textasciitilde{}\textasciitilde{}\textasciitilde{} \_ t:tetrahedra (- i \_i j\_i(t)) j\_1(t)/j\_2(t)/j\_3(t)/j\_4(t)/j\_5(t)/j\_6(t)/ .
\end{document}
